\documentclass[11pt]{article}
\usepackage{icmlko}
\begin{document}

\title{빌려 쓰는 도메인은 인질이다 --- 자기 행이 한 줄도 안 바뀌었는데
$0.31$ 을 잃는다}
\author{월드모델 실험실}
\date{노트 452 · 2026년 8월 2일}
\maketitle

\begin{abstract}
노트 451 이 최근성 이득 $+0.0057$ 을 판 전체 평균으로 냈다. 도메인마다
다를 것 같아 갈라 봤다. 등록한 문장은 \emph{``가장 큰 도메인과 가장 작은
도메인의 차가 뽑기 SE 의 네 배를 넘는다''} 였다 --- \textbf{$3.9$배로
아슬하게 못 넘었다}. 노트 451 이 방금 만든 규약대로 \textbf{반증으로
적는다.}

그런데 표를 읽다가 다른 것이 보였다.

\begin{center}
\begin{tabular}{lrrr}
\hline
도메인 & 차 & SE & 자기 행이 바뀐 비율 \\
\hline
애니 & $+0.0349$ & $0.0046$ & $42\%$ \\
세계애니 & $+0.0339$ & $0.0044$ & $70\%$ \\
만화 & $-0.0250$ & $0.0065$ & $75\%$ \\
도서 & $-0.0372$ & $0.0179$ & $58\%$ \\
\textbf{팝업} & $\mathbf{-0.3089}$ & $0.0122$ & $\mathbf{0\%}$ \\
\hline
\end{tabular}
\end{center}

\textbf{팝업은 자기 학습 행이 한 줄도 안 바뀌었다.} 학습 $24$행이 전부
$2020$년 이후라 ``최근 $N$''과 ``무작위 $N$''이 \emph{같은 행}이다. 그런데
$-0.3089$ 를 잃는다. \textbf{잃은 이유가 전부 남에게 있다.}
\end{abstract}

\begin{figure}[t]
\centering
\includegraphics[width=\textwidth]{figs/f1.pdf}
\caption{\textbf{가} --- 도메인별 최근성 이득. 초록이 자기 SE 두 배를
넘는 이득(다섯), 회색이 못 정함, 붉은색이 \emph{자기 행이 하나도 안
바뀐} 도메인. \textbf{나} --- 학습 행이 적을수록 크게 흔들린다.}
\end{figure}

\section{노트 445 가 예고한 것}

노트 445 는 시장팝업의 같은 시기 천장이 $-0.0901$ 이라 \emph{``가진 게
전부 빌린 것''}이라고 적었다. 그때는 관찰이었다. 이번에 그것이
\textbf{측정}이 된다.

자기 행을 하나도 안 바꾼 도메인이 둘이다 --- 팝업(학습 $24$행)과
시장팝업(학습 $123$행). 둘 다 손해를 보는데 \textbf{크기가 스물다섯 배
차이}다($-0.3089$ 대 $-0.0122$). 학습 행이 적을수록, 즉 \emph{남에게
더 많이 빌릴수록} 크게 흔들린다.

\textbf{빌리는 쪽은 빌려주는 쪽의 구성에 인질로 잡혀 있다.} 판이
``최근 행만 쓰자''로 바뀌면 팝업은 자기가 아무것도 안 했는데 점수의
$70\%$ 를 잃는다.

\section{이것이 규약을 하나 바꾼다}

지금 관문은 \textbf{판 전체}와 \textbf{도메인별 부호}를 본다. 그런데
도메인별 부호는 \emph{그 도메인의 자료가 바뀌었는지}를 안 본다. 팝업처럼
자기 행이 안 바뀐 도메인의 $-0.3089$ 는 ``팝업에 나쁜 변경''이 아니라
\textbf{``팝업이 남에게 얼마나 의존하는지''}를 재고 있다.

\begin{quote}
\textbf{조항} --- 학습 자료의 구성을 바꾸는 변경은 도메인별 결과를 적을
때 \textbf{그 도메인의 학습 행이 얼마나 바뀌었는지를 같이 적는다.}
$0\%$ 바뀐 도메인의 점수 변화는 그 도메인에 대한 판정이 아니라
\emph{의존도의 측정}이다.
\end{quote}

\section{등록한 문장은 왜 못 넘었나}

$3.9$배는 아쉬운 숫자지만 규약대로 반증이다. 다만 \emph{왜} 못 넘었는지는
적어 둔다 --- 내가 고른 통계(최대$-$최소)가 \textbf{양 끝의 SE 를 둘 다
짊어진다}. 최대인 애니(SE $0.0046$)와 최소인 도서(SE $0.0179$)를 합치면
$0.0185$ 인데, 도서는 학습 $80$행이라 SE 가 크다. \textbf{끝값 두 개로
이질성을 재면 제일 안 미더운 도메인이 판정을 쥔다.}

다르게 읽으면 그림은 분명하다 --- \textbf{다섯 도메인이 자기 SE 의 두 배를
넘는 이득}(애니 · 세계애니 · 게임 · 모바일 · 펀딩)이고 \textbf{셋이 그만큼
손해}(도서 · 만화 · 팝업)다. 하지만 그 통계는 \textbf{등록 안 했으므로}
이 노트의 결론이 아니라 다음 노트의 등록거리로 남긴다.

\section*{남은 것}

\begin{itemize}
\item \textbf{다음에 등록할 것} --- ``자기 SE 두 배를 넘는 도메인이
양쪽에 다 있다''를 미리 적고 다시 잰다. 끝값이 아니라 개수로 재는 것이
힘이 세다.
\item \textbf{만화가 왜 옛 행을 원하나}($-0.0250$). 제일 큰 도메인이고
최근 비율이 제일 낮다($25\%$). 형식이 안 변해서인가, 아니면 옛 행이
많아서 최근만 쓰면 손해가 큰 것인가.
\item \textbf{의존도를 직접 재기.} 팝업 $-0.3089$ 는 우연히 나온
의존도 측정이다. 도메인마다 \emph{남의 자료를 흔들었을 때 얼마나
흔들리는지}를 대놓고 재면 노트 445 의 ``빌린 것''을 수치로 만들 수 있다.
\end{itemize}

\end{document}
